\documentclass[10pt,a4paper]{amsart}
\usepackage[margin=19mm]{geometry}
\usepackage{amsmath,amssymb,mathtools}
\usepackage[scr=rsfs,cal=euler]{mathalfa}
\usepackage{xfrac}
\usepackage[unicode,hidelinks,bookmarks=false]{hyperref}
\newcommand{\ie}{i.e.\ }
\newcommand{\cf}{cf.\ }
\newcommand{\resp}{respectively\ }
\newcommand{\Subsection}{\subsection}
\providecommand{\nobreakdash}{\nobreak\hbox{-}}
\setlength{\emergencystretch}{2em}
\multlinegap0pt
\usepackage[shortlabels]{enumitem}
\usepackage{amssymb}
\usepackage{mathtools}
\usepackage{xcolor}
\newtheorem{lemma}{Lemma}
\usepackage{cleveref}
\crefname{lemma}{Lemma}{Lemmas}
\newcommand{\ZZ}{\mathbb Z}
\title{Sparse $k$-AP Covering Sets and the Arithmetic Kakeya Conjecture}
\author{Pitchayut Saengrungkongka}
\date{Source: arXiv:2609.02041v1 (CC BY 4.0)}
\begin{document}
\maketitle

\noindent\emph{Source and licence.} Sparse $k$-AP Covering Sets and the Arithmetic Kakeya Conjecture. Version arXiv:2609.02041v1 (CC BY 4.0), distributed by arXiv under the Creative Commons Attribution 4.0 International licence, http://creativecommons.org/licenses/by/4.0/, as recorded in the arXiv licence metadata for this exact version and shown on the abstract page https://arxiv.org/abs/2609.02041v1. Full licence notice: \texttt{licenses/arxiv-2609.02041-CC-BY-4.0.txt}.\par\smallskip

\noindent\textit{Editorial scope.} Counted unit: the article's lemma lem:sampling (source0.tex lines 492--503). The article's complete original proof of it is delivered with it (source0.tex lines 504--594, one \texttt{proof} environment of 91 lines). No mathematics is written, repaired or re-derived: the statement and the proof are the article's own lines, in the article's own notation, and no retained line is substituted or re-flowed.  Retained uncounted as the source-local context the proof needs: lines 405--416; lines 463--474.  Nothing was omitted from any retained span, so the package carries no omission record.  No retained line contains a citation, so the wrapper carries no bibliography and no bibliography entry.  No retained line contains a figure environment, an image-inclusion command or drawing code, so no image is delivered and the package declares no asset dependency.  Editorial formatting only, in addition: an AMS \texttt{amsart} preamble that restates the article's own declarations and macros used by the retained lines, namely the environments \texttt{lemma} on separate counters, together with the standard packages those lines need.  The retained environments print in this document as Lemma 1 in order of appearance, whereas the article prints the same items with the numbering of its own full document.  The complete native source of the article as distributed in the arXiv e-print for this exact version is bundled unchanged as \texttt{sources/209166/source0.tex}.
\begin{lemma}
\label{lem:crt}
There are infinitely many positive integers $Q$ 
for each of which there exists 
a function $f : \ZZ/Q\ZZ \to (0,\infty)$ such that 
\begin{enumerate}[label=(\alph*)]
\item $0<f(x) < Q^{\frac 1{100k}}$ for all $x\in \ZZ/Q\ZZ$.
\item $\mathbb E_x[f(x)] = 1$.
\item for any $a\in \ZZ/Q\ZZ$, we have
$$\mathbb E_x[f(a-3x)f(a-4x)f(a-5x)] > 100^{k^2}\log Q.$$
\end{enumerate}
\end{lemma}

\begin{lemma}
\label{lem:smoothing}
For each $Q$ in \Cref{lem:crt},
there exists a function $g : \ZZ/Q\ZZ\to (0,\infty)$
such that 
\begin{enumerate}[label=(\alph*)]
\item $0<g(x) < Q^{\frac 1{100k}}$ for all $x\in \ZZ/Q\ZZ$.
\item $\mathbb E_x[g(x)] = 1$.
\item for any $a\in \ZZ/Q\ZZ$, we have
$$\mathbb E_x\left[\prod_{j=1}^{k-1} g(a-jx)\right] > 50^{k^2}\log Q.$$
\end{enumerate}
\end{lemma}

\begin{lemma}
\label{lem:sampling}
There exist a positive integer $N$ 
and a set $B\subseteq \{0,1,\dots,N-1\}$
such that 
\begin{enumerate}[label=(\alph*)]
\item $|B| < \frac 12 N^{\frac{k-2}{k-1}}$ and 
\item for any $a\in \{0,1,\dots,N-1\}$, there exists 
a nonnegative integer $d$ such that
$a-d$, $a-2d$, \dots, $a-(k-1)d$ are all in $B$.
\end{enumerate}
\end{lemma}

\begin{proof}
We take $Q$ and $g$ as in \Cref{lem:smoothing}.
Let $N=10k^2Q^2$. 
We construct the set $B$ as follows:
\begin{itemize}
    \item include each $x\in [0,10k^2Q)$;
    \item include each $x\in [10k^2Q, 20k^2Q)$ 
    with probability $\frac 1{100k^2} Q^{-\frac 1{k-1}}g(x)$;
    \item include each $x\in [20k^2Q, 30k^2Q)$
    with probability $\frac 1{100k^2} (2Q)^{-\frac 1{k-1}} g(x)$;
    \item include each $x\in [30k^2Q, 40k^2Q)$
    with probability $\frac 1{100k^2} (3Q)^{-\frac 1{k-1}} g(x)$;
    \item \hspace{2cm} \vdots 
    \item include each $x\in [10(Q-1)k^2Q, 10k^2Q^2)$
    with probability $\frac 1{100k^2} ((Q-1)Q)^{-\frac 1{k-1}} g(x)$.
\end{itemize}
In particular, we partition $\{0,1,\dots,N-1\}$ 
into $Q$ blocks, each of which has length $10k^2Q$.
Note that 
\begin{align*}
\mathbb E[|B|] &= 10k^2Q + 
\Big(1+2^{-\frac 1{k-1}}+3^{-\frac 1{k-1}} + \dots + (Q-1)^{-\frac 1{k-1}}\Big)
\cdot \left(\frac 1{100k^2}Q^{-\frac 1{k-1}}\right) \sum_{x=0}^{10k^2Q-1} g(x)\\
&< 20k^2Q  + \left(\frac{k-1}{k-2}\cdot  Q^{\frac{k-2}{k-1}}\right)
\cdot \left(\frac 1{100k^2}  Q^{-\frac 1{k-1}}\right) \cdot 
\left(10k^2Q \cdot \mathbb E_x[g(x)]\right) \\
&< 20k^2 Q + \frac 15 Q^{2\cdot \frac{k-2}{k-1}} \\
&< \frac 14 Q^{2\cdot \frac{k-2}{k-1}},
\end{align*}
if $Q$ is sufficiently large.
Thus, by Markov's inequality, we deduce that 
$$\mathbb P\left(|B| > \tfrac 12 N^{\frac{k-2}{k-1}}\right) < 0.1.$$

Next, we look at the second condition.
Let $a\in\{0,1,\dots,N-1\}$, and take $n$ such that
$a\in [10k^2nQ, 10k^2(n+1)Q)$.
If $n=0$, then automatically we have $a\in B$,
so we may pick $d=0$.

Henceforth, we assume that $n\geq 1$.
We commit to picking $x$ from the interval $[knQ, (k+1)nQ]$. 
We claim that for any distinct $x_1, x_2\in [knQ, (k+1)nQ]$,
the sets $\{a-x_1, a-2x_1,\dots,a-(k-1)x_1\}$
and $\{a-x_2,a-2x_2,\dots,a-(k-1)x_2\}$ are disjoint.
Suppose that they are not disjoint. Then there exists 
$u,v\in\{1,2,\dots,k-1\}$ such that 
$ux_1 = vx_2$, so $\tfrac{x_1}{x_2} = \tfrac uv$.
However, $\tfrac{x_1}{x_2}$ is between $\tfrac k{k+1}$
and $\tfrac{k+1}k$, which cannot be equal to 
$\tfrac uv$ for $u,v\in\{1,2,\dots,k-1\}$
(unless $u=v$).
Therefore, $u=v$. From $ux_1=vx_2$, we then conclude that 
$x_1=x_2$, which is a contradiction.

Thus, we deduce that the events of the form 
$\{a-x, a-2x,\dots,a-(k-1)x\} \subseteq B$ 
are independent across all $x$.
Furthermore, since $a-jx < 10k^2(n+1)Q$
(so it is in at most the $(n+1)$-th block)
and the prefactors assigned to each block is decreasing,
we get that the probability that $a-jx\in B$ is at least 
$\frac 1{100k^2}(nQ)^{-\frac{1}{k-1}}g(a-jx)$. 

Therefore, the probability that $a$ does not 
satisfy the second condition is
\begin{align*}
\mathbb P(a\text{ fails}) 
&= \prod_{x=knQ}^{(k+1)nQ} \Big(1 - \prod_{j=1}^{k-1} 
\mathbb P[a-jx\in B]\Big) \\
&\leq \prod_{x=knQ}^{(k+1)nQ} \Big(1 - \prod_{j=1}^{k-1} 
\left(\frac 1{100k^2}\cdot 
(nQ)^{-\frac 1{k-1}} g(a-jx)\right)\Big) \\
&= \prod_{x=knQ}^{(k+1)nQ} \Big(1 - 
\frac 1{(100k^2)^{k-1} nQ} \prod_{j=1}^{k-1} g(a-jx)\Big) \\
&< \prod_{x=knQ}^{(k+1)nQ} \Big(1 - 
\frac 1{10^{k^2} nQ} \prod_{j=1}^{k-1} g(a-jx)\Big) \\
&\leq \prod_{x=knQ}^{(k+1)nQ}  \exp\left(-\frac 1{10^{k^2}nQ}
\prod_{j=1}^{k-1} g(a-jx) \right) \tag{using $\exp(-x)\geq 1-x$}\\
&\leq \exp\left(-\frac 1{10^{k^2}nQ} \sum_{x=knQ}^{(k+1)nQ}
\prod_{j=1}^{k-1} g(a-jx) \right)\\
&\leq \exp\Big(-\frac 1{10^{k^2}} 
\mathbb E_x \left[\prod_{j=1}^{k-1} g(a-jx)\right]\Big) \\
&\leq \exp\Big(-100\log Q\Big) < \frac 1{Q^{100}},
\end{align*}
so summing across all $a\in \{0,1,\dots,N-1\}$ 
and using the union bound gives 
that the probability that the second condition fails is at most $\frac 12$.
Thus, by the union bound, there is a nonzero probability
that $B$ satisfies both conditions.
Therefore, such a $B$ exists.
\end{proof}

\end{document}
